\begin{tikzpicture}[x=1.05cm,y=0.5cm,
  lab/.style={font=\scriptsize,align=center},
  dot/.style={fill=black,draw=none}]
\def\colw{0.85}
% ---- five stage columns: graphene sheets (lines) ----
\foreach \x/\name in {0/{stage 4}, 1.4/{stage 3}, 2.8/{stage 2L}, 4.2/{stage 2}, 5.6/{stage 1}}{
  \foreach \k in {0,...,6}{\draw[semithick] (\x,0.6*\k)--(\x+\colw,0.6*\k);}
  \node[lab,below] at (\x+\colw/2,-0.25) {\name};
}
% ---- occupied galleries: discrete Li dots, ordered rows for crystalline
%      stages, jittered rows for the liquid-like stage 2L ----
% stage 4: gallery 1 and 5 filled, ordered
\foreach \k in {1,5}{\foreach \dx in {0.12,0.30,0.48,0.66,0.84}{\fill[dot] (0+\dx,0.6*\k-0.30) circle (0.38pt);}}
% stage 3: gallery 1 and 4 filled, ordered
\foreach \k in {1,4}{\foreach \dx in {0.12,0.30,0.48,0.66,0.84}{\fill[dot] (1.4+\dx,0.6*\k-0.30) circle (0.38pt);}}
% stage 2L: gallery 1,3,5 filled, DISORDERED (in-plane order lost, liquid-like)
\foreach \k/\ya in {1/0.02,3/-0.03,5/0.015}{
  \foreach \dx/\dy in {0.08/0.01,0.24/-0.02,0.40/0.03,0.55/-0.015,0.70/0.02,0.84/-0.01}{
    \fill[dot] (2.8+\dx,0.6*\k-0.30+\ya+\dy) circle (0.34pt);
  }
}
% stage 2: gallery 1,3,5 filled, ordered dense
\foreach \k in {1,3,5}{\foreach \dx in {0.10,0.26,0.42,0.58,0.74}{\fill[dot] (4.2+\dx,0.6*\k-0.30) circle (0.40pt);}}
% stage 1: all galleries filled, ordered dense (LiC6, full occupancy)
\foreach \k in {1,...,6}{\foreach \dx in {0.10,0.26,0.42,0.58,0.74}{\fill[dot] (5.6+\dx,0.6*\k-0.30) circle (0.40pt);}}
% ---- direction arrows ----
\draw[-{Latex[length=1.8mm]},thick] (0,4.05)--(6.45,4.05) node[pos=0.5,above,font=\scriptsize] {lithiation (charge): $4\to3\to2\mathrm L\to2\to1$};
\draw[{Latex[length=1.8mm]}-,thick] (0,-1.05)--(6.45,-1.05) node[pos=0.5,below,font=\scriptsize] {delithiation (discharge): main direction, $\theta_j:1\to0$};
% ---- U_j / peak-position axis: ties each stage transition to its
%      literature initial value and the dQ/dV peak it leaves (tab:staging) ----
\draw[-{Latex[length=1.5mm]}] (-0.3,-2.05) -- (6.75,-2.05) node[right,font=\scriptsize] {};
\foreach \x/\uval in {1.125/0.210, 2.525/0.140, 3.925/0.120, 5.325/0.085}{
  \draw[semithick] (\x,-2.05) -- (\x,-1.95);
  \draw[thick] (\x-0.16,-1.95) -- (\x,-1.62) -- (\x+0.16,-1.95);
  \node[font=\scriptsize] at (\x,-2.28) {$U{=}\uval$ V};
}
\node[font=\scriptsize,align=center,anchor=west] at (6.85,-2.05) {$U_j$ initial\\(tab.~\ref{tab:staging})};
\end{tikzpicture}
\caption{흑연 staging 의 갤러리 채움 도식. 수평선이 그래핀 층, 점 무늬가 층간 갤러리에 든 리튬이다 —
결정성 stage(4/3/2/1)는 갤러리 내 리튬을 정렬된 열로, 면내 질서가 풀린 stage $2\mathrm L$ 은 같은 갤러리를
흐트러진 점 배치로 그려 액체상(liquid-like) 임을 구분한다. 채울수록 stage 번호가 줄어 stage
1($\mathrm{LiC_6}$, 완전 채움)에 닿고, 각 stage 전이가 $\dd Q/\dd V$ 에 peak 하나를 남긴다. 아래 뾰족 기호와
숫자는 코드의 전이 리스트 4건이 갖는 초기값 $U$(표~\ref{tab:staging}, 문헌 기반)를 그 전이가 놓인 칼럼
경계에 정렬한 것이다 — stage 번호가 줄수록(우측일수록) $U$ 도 줄어, 방전(탈리튬화, 하단 화살표) 진행 중
전위가 낮은 전이($2\to1$, $0.085$ V)부터 순서대로 지난다. 본 문건 주 전개 방향은 방전이다.}
\label{fig:staging}
